
\documentclass[12pt]{article}
\usepackage{amsmath}
\usepackage{amsfonts}
\usepackage{amssymb}
\usepackage{geometry}
\geometry{a4paper, margin=0.8in}
\pagenumbering{gobble}

\title{Homework assignment for 02YELMA \\ Chapter: Electrostatics}
\begin{document}
\date{}
\maketitle
\vspace{-0.8in}

\begin{enumerate}
    \item Four equal charges, $+q$, are located at the corners of a square with side length $a$. What is the net electric field in the center of the square?  If we replace one of the charges with $-q$, what would be the electric field and what would be the net force on a test charge $Q$ in the center?
    
    \item Find the electric field a distance $z$ above the center of a circular loop of radius $r$, which carries a uniform line charge $\lambda$.
    
    \item  Find the electric field inside and outside of a charged sphere with charge density $\rho=kr$, where $k$ is a constant and $r$ is the distance from the center of the sphere.
    
    \item Depending on the result of Q3, calculate the potential inside and outside of the sphere.

    \item If you consider the situation in Q1 again (four equal charges, $+q$, are situated at the corners of a square), how would the field lines look?

    \item Consider a uniform disc-shaped surface charge $\sigma$ with radius $R$. Find the potential at a distance $z$ above its center. Calculate the electric field using the potential. (Hint: Consider the substitution $u=r^2+z^2$)
    
    \item  If we assume that the electron has a uniform surface charge distribution in the form of a spherical shell, and we calculate its rest energy using the formula $E=mc^2$, what would be its radius \textit{in meters}? Find relevant quantities on the Internet such as $m$, $c$, etc., ensuring that they are compatible in terms of units.
    
    \item A metal sphere of radius $R$ carries a total charge $Q$. What is the force of repulsion between the northern and southern hemispheres?
    
    \item Suppose that the plates of a parallel plate capacitor move closer together by an infinitesimal distance $\epsilon$, as a result of their mutual interaction. What is the amount of work done by electrostatic forces, in terms of the electric field $E$ and the area of the plates $A$? (Hint: Use $P=\frac{\epsilon_0}{2}E^2$)

    \item Find the capacitance \textit{per unit length} of two cylindrical coaxial metal tubes, of radii $a$ and $b$.

    \item Four particles are placed in the following positions: $-2q \rightarrow (0, -a, 0)$, $-2q \rightarrow (0, a, 0)$, $q \rightarrow (0, 0, -a)$, $3q \rightarrow (0, 0, a)$. Find a simple approximate formula for the potential that is valid at points far from the origin.

    \item Two charges are placed in two different arrangements: (a) $3q$ at $(0,0,a)$, and $-q$ at $(0,0,0)$, and (b) $3q$ at $(0,0,0)$, and $-q$ at $(0,0,-a)$. For each arrangement, determine the monopole moment (in other words, the monopole contribution), the dipole moment, and the approximate potential (in spherical coordinates) using these moments.
    
    \item A pure dipole $p$ is located at the origin, oriented in the $-z$ direction. What is the force experienced by a point charge located at the coordinates $(a,0,0)$?

    \item Two pure electric dipoles $\vec{p}_1$ and $\vec{p}_2$ are a distance $r$ apart along the $x$-axis. What is the torque experienced by $\vec{p}_2$ due to the presence of $\vec{p}_1$? ($\vec{p}_1$ and $\vec{p}_2$ are oriented along the $\hat{z}$ and $\hat{x}$ directions, respectively.)

    \item A sphere of radius $R$ has a polarization $\vec{P}(\vec{r})=k\vec{r}$ where $k$ is a constant and $\vec{r}$ is the vector from the center. Calculate the bound charges $\sigma_b$ and $\rho_b$, and find the field inside and outside the sphere.

\end{enumerate}

\underline{Given for 13 and 14}: The electric field created by a pure dipole:
\begin{equation*}
    \vec{E}_{dip}(r,\theta) = \frac{p}{4\pi\epsilon_0 r^3}\left[2\cos{\left(\theta\right)}\hat{r} + \sin{\left(\theta\right)}\hat{\theta}\right]
\end{equation*}


\end{document}
